\documentclass[11pt,a4paper]{article}

\usepackage[T1]{fontenc}
\usepackage[utf8]{inputenc}
\usepackage[polish]{babel}
\usepackage{amsmath, amssymb, amsthm}
\usepackage{mathtools}
\usepackage{geometry}
\geometry{margin=2.5cm}

\theoremstyle{definition}
\newtheorem{definition}{Definicja}
\newtheorem{notation}{Notacja}
\newtheorem{lemma}{Lemat}
\newtheorem{observation}{Obserwacja}
\newtheorem{proposition}{Propozycja}
\newtheorem{theorem}{Twierdzenie}

\newcommand{\len}[1]{\lVert #1 \rVert}
\newcommand{\dist}[3]{d_{#1}(#2,#3)}

\title{\textbf{Zbiory semi-rozcinające w grafach}}
\author{}
\date{}

\begin{document}
\maketitle

\section{Notacja i definicje}

W całym tekście $G=(V,E)$ oznacza skończony graf nieskierowany. Nie zakładamy spójności grafu. Jeżeli między
dwiema różnymi składowymi spójności nie istnieje żadna ścieżka, to odpowiadającą odległość definiujemy jako~$\infty$.

\begin{definition}[Ścieżki i odległość]\label{def:paths-distance}
Spacerem z~$u$ do~$w$ nazywamy ciąg wierzchołków $(v_0,v_1,\dots,v_k)$ taki, że $v_0=u$, $v_k=w$ oraz
$v_{i-1}v_i\in E$ dla $i=1,\dots,k$. Długość takiego spaceru wynosi $\len{(v_0,\dots,v_k)}\coloneqq k$.
Ścieżką nazywamy spacer bez powtórzeń wierzchołków.

Odległość $\dist{G}{u}{w}$ definiujemy jako najmniejszą długość spaceru z~$u$ do~$w$ w~$G$, o~ile taki spacer
istnieje; w~przeciwnym razie $\dist{G}{u}{w}\coloneqq\infty$.
\end{definition}

\begin{notation}[Wierzchołki i krawędzie ścieżki]\label{not:path-sets}
Jeżeli $P=(v_0,\dots,v_k)$ jest spacerem w~$G$, to oznaczamy
\[
  V(P)\coloneqq \{v_0,\dots,v_k\},
  \qquad
  E(P)\coloneqq \{v_{i-1}v_i:\ i=1,\dots,k\}.
\]
\end{notation}

\begin{notation}[Usuwanie krawędzi]\label{not:edge-removal}
Jeżeli $F\subseteq E$, to $G\setminus F$ oznacza graf $(V,E\setminus F)$.
\end{notation}

\begin{notation}[Końce krawędzi]\label{not:VF}
Dla $F\subseteq E$ przez $V[F]$ oznaczamy zbiór wierzchołków incydentnych z~krawędziami z~$F$, tj.
\[
  V[F]\coloneqq \{v\in V:\ \exists\, e\in F\ \text{takie, że $e$ jest incydentna z~$v$}\}.
\]
\end{notation}

\begin{notation}[Sąsiedztwo]\label{not:neighborhood}
Jeżeli $v\in V$, to $N(v)\coloneqq\{w\in V:\ vw\in E\}$ oraz $N[v]\coloneqq N(v)\cup\{v\}$.
\end{notation}

\begin{notation}[Ściąganie zbioru wierzchołków]\label{not:contraction}
Niech $W\subseteq V$ i~niech $w^\star\notin V$ będzie nowym wierzchołkiem. Definiujemy odwzorowanie
$\pi_W\colon V\to (V\setminus W)\cup\{w^\star\}$ dane wzorem
\[
  \pi_W(v)\coloneqq
  \begin{cases}
    w^\star & \text{gdy } v\in W,\\
    v & \text{gdy } v\notin W.
  \end{cases}
\]
Graf $G/W$ definiujemy jako $G/W\coloneqq (V',E')$, gdzie
\[
  V'\coloneqq (V\setminus W)\cup\{w^\star\},
  \qquad
  E'\coloneqq \{\pi_W(u)\pi_W(v):\ uv\in E,\ \pi_W(u)\neq \pi_W(v)\}.
\]
Krawędzie o obu końcach w~$W$ znikają (po ściągnięciu dawałyby pętle). Dla krawędzi $uv\in E$ z~co najwyżej
jednym końcem w~$W$ jej obraz w~$G/W$ oznaczamy przez $uv/W\coloneqq \pi_W(u)\pi_W(v)\in E'$.
Dla zbioru krawędzi $F\subseteq E$ definiujemy jego obraz
\[
  F/W\coloneqq \{\, e/W : e\in F\ \text{i $e/W$ jest zdefiniowane}\,\}\subseteq E'.
\]
\end{notation}

\begin{definition}[Zbiór semi-rozcinający]\label{def:semi-cut}
Niech $G=(V,E)$ będzie grafem. Mówimy, że $F\subseteq E$ jest zbiorem semi-rozcinającym dla $T\subseteq V$ w~$G$,
jeżeli dla każdej pary różnych wierzchołków $u,w\in T$ zachodzi
\[
  \dist{G}{u}{w} \;<\; \dist{G\setminus F}{u}{w}.
\]
\end{definition}

\section{Lematy i propozycje pomocnicze}

\begin{lemma}\label{lem:walk-to-path}
Jeżeli istnieje spacer z~$u$ do~$w$ w~$G$, to istnieje również ścieżka z~$u$ do~$w$ w~$G$ o~długości nie większej niż
długość tego spaceru. W~szczególności, jeżeli $\dist{G}{u}{w}<\infty$, to istnieje ścieżka z~$u$ do~$w$ o~długości
dokładnie $\dist{G}{u}{w}$.
\end{lemma}

\begin{proof}
Weźmy dowolny spacer $P=(v_0,\dots,v_k)$ z~$u$ do~$w$. Jeżeli w~$P$ występuje powtórzenie wierzchołka, to istnieją
indeksy $0\le i<j\le k$ takie, że $v_i=v_j$. Usuwając z~ciągu wierzchołki $(v_{i+1},\dots,v_j)$ otrzymujemy spacer z~$u$
do~$w$ o~mniejszej długości. Powtarzając tę operację skończenie wiele razy, dostajemy spacer bez powtórzeń wierzchołków,
czyli ścieżkę, o~długości nie większej od $\len{P}$. Druga część wynika z~definicji odległości jako minimum po długościach
spacerów.
\end{proof}

\begin{lemma}\label{lem:subpath}
Niech $P=(v_0,\dots,v_k)$ będzie ścieżką o~długości $\dist{G}{u}{w}$. Niech $x=v_i$ i~$y=v_j$ dla pewnych $0\le i\le j\le k$
i~niech $P[x,y]\coloneqq (v_i,\dots,v_j)$ będzie odpowiadającą podścieżką $P$ łączącą $x$ i~$y$. Wówczas $\len{P[x,y]}=\dist{G}{x}{y}$.
\end{lemma}

\begin{proof}
Niech $\len{P[x,y]}=L$. Gdyby istniała ścieżka $Q$ z~$x$ do~$y$ o~długości mniejszej niż $L$, to podmieniając w~$P$ odcinek
$P[x,y]$ na $Q$ otrzymalibyśmy ścieżkę z~$u$ do~$w$ krótszą niż $\dist{G}{u}{w}$, co jest sprzeczne z~minimalnością $P$.
\end{proof}

\begin{observation}\label{obs:shortest-paths}
Niech $G=(V,E)$ będzie grafem oraz niech $F\subseteq E$ i~$T\subseteq V$. Wówczas $F$ jest zbiorem semi-rozcinającym dla~$T$
w~$G$ wtedy i~tylko wtedy, gdy dla każdych różnych $u,w\in T$ spełnione są łącznie następujące warunki:
$\dist{G}{u}{w}<\infty$ oraz każda ścieżka $P$ z~$u$ do~$w$ o~długości $\len{P}=\dist{G}{u}{w}$ zawiera co najmniej jedną
krawędź ze zbioru $F$.
\end{observation}

\begin{proof}
Ustalmy różne $u,w\in T$.

(\,$\Rightarrow$\,) Z~Definicji~\ref{def:semi-cut} mamy $\dist{G}{u}{w}<\dist{G\setminus F}{u}{w}\le\infty$, a~zatem
$\dist{G}{u}{w}<\infty$. Niech $P$ będzie ścieżką z~$u$ do~$w$ o~długości $\dist{G}{u}{w}$.
Gdyby $P$ nie zawierała żadnej krawędzi z~$F$, to byłaby również ścieżką w~$G\setminus F$ o~tej samej długości, co dawałoby
$\dist{G\setminus F}{u}{w}\le \dist{G}{u}{w}$, sprzecznie z~Definicją~\ref{def:semi-cut}.

(\,$\Leftarrow$\,) Załóżmy, że $\dist{G}{u}{w}<\infty$ oraz każda ścieżka $P$ z~$u$ do~$w$ o~długości $\dist{G}{u}{w}$
zawiera krawędź z~$F$. Rozważmy dowolną ścieżkę $Q$ z~$u$ do~$w$ w~$G\setminus F$. Jest ona jednocześnie ścieżką w~$G$ i~z
definicji nie zawiera krawędzi z~$F$, więc nie może mieć długości $\dist{G}{u}{w}$. Otrzymujemy $\len{Q}>\dist{G}{u}{w}$ dla
każdej takiej $Q$, co oznacza $\dist{G\setminus F}{u}{w}>\dist{G}{u}{w}$. Ponieważ $u,w$ były dowolne, $F$ spełnia
Definicję~\ref{def:semi-cut}.
\end{proof}

\begin{lemma}\label{lem:projection-walk}
Niech $G=(V,E)$ będzie grafem, $W\subseteq V$, a~$G/W$ zdefiniowany jak w~Notacji~\ref{not:contraction}.
Jeżeli $(v_0,\dots,v_k)$ jest spacerem w~$G$, to ciąg $(\pi_W(v_0),\dots,\pi_W(v_k))$ po usunięciu kolejnych powtórzeń
wierzchołków jest spacerem w~$G/W$ o~długości nie większej niż $k$.
\end{lemma}

\begin{proof}
Jeżeli $v_{i-1}v_i\in E$ oraz $\pi_W(v_{i-1})\neq \pi_W(v_i)$, to z~definicji $E'$ mamy $\pi_W(v_{i-1})\pi_W(v_i)\in E'$.
Jeżeli natomiast $\pi_W(v_{i-1})=\pi_W(v_i)$, to dany krok spaceru znika po usunięciu kolejnych powtórzeń. W~ten sposób
otrzymujemy spacer w~$G/W$; każdy usunięty krok zmniejsza długość, więc długość wyniku jest nie większa niż $k$.
\end{proof}

\begin{proposition}\label{prop:contraction}
Niech $G=(V,E)$ będzie grafem, $T\subseteq V$, $F\subseteq E$ i~niech $F$ będzie zbiorem semi-rozcinającym dla~$T$ w~$G$.
Niech $t\in T\setminus V[F]$ oraz niech $t'\notin V$ będzie nowym wierzchołkiem. Rozważmy graf $G'\coloneqq G/N[t]$, gdzie
$N[t]$ zostało ściągnięte do wierzchołka~$t'$, zbiór wierzchołków $T'\coloneqq (T\setminus\{t\})\cup\{t'\}$ oraz zbiór krawędzi
$F'\coloneqq F/N[t]\subseteq E(G')$. Wówczas $F'$ jest zbiorem semi-rozcinającym dla~$T'$ w~$G'$.
\end{proposition}

\begin{proof}
Najpierw pokażemy, że żaden wierzchołek z~$T\setminus\{t\}$ nie leży w~$N(t)$. Weźmy $u\in T\setminus\{t\}$.
Jeżeli $u\in N(t)$, to $\dist{G}{t}{u}=1$, a~każda ścieżka z~$t$ do~$u$ o~długości $1$ składa się z jednej krawędzi $tu$.
Z~Obserwacji~\ref{obs:shortest-paths} wynikałoby $tu\in F$, czyli $t\in V[F]$, sprzecznie z~założeniem.
Zatem $T\setminus\{t\}\subseteq V\setminus N[t]$ i~w szczególności $T'\subseteq V(G')$.

Ustalmy różne $u,w\in T'$ i~niech $P'$ będzie ścieżką z~$u$ do~$w$ w~$G'$ o~długości $\len{P'}=\dist{G'}{u}{w}$.
Pokażemy, że $P'$ zawiera krawędź ze zbioru $F'$.

\medskip\noindent
\textbf{Przypadek 1: $t'\notin V(P')$.}
Wtedy $P'$ jest również ścieżką w~$G$. Gdyby w~$G$ istniała ścieżka $Q$ z~$u$ do~$w$ o~długości mniejszej niż $\len{P'}$,
to na mocy Lematu~\ref{lem:projection-walk} jej obraz w~$G'$ byłby spacerem z~$u$ do~$w$ o~długości nie większej niż
$\len{Q}<\len{P'}$, co przeczyłoby minimalności $P'$. Zatem $\len{P'}=\dist{G}{u}{w}$.
Ponieważ $u,w\in T\setminus\{t\}$, z~Obserwacji~\ref{obs:shortest-paths} wynika, że $P'$ zawiera pewną krawędź $e\in F$.
Krawędź $e$ ma oba końce poza $N[t]$, więc $e/N[t]=e$ i~w szczególności $e\in F'$. Stąd $E(P')\cap F'\neq\emptyset$.

\medskip\noindent
\textbf{Przypadek 2: $t'\in V(P')$.}
Najpierw załóżmy, że $w=t'$ (przypadek $u=t'$ jest symetryczny). Niech $P'$ ma postać $(t',x_1,\dots,x_k)$, gdzie $x_k=u$.
Z~definicji ściągania wynika, że krawędź $t'x_1$ w~$G'$ pochodzi od pewnej krawędzi $vx_1$ w~$G$ z~wierzchołkiem $v\in N[t]$.
Ponieważ $x_1\notin N[t]$, w~istocie $v\in N(t)$, a~więc $tv\in E(G)$.
Rozważmy ścieżkę $P$ w~$G$ daną ciągiem wierzchołków $(t,v,x_1,\dots,x_k)$. Mamy $\len{P}=\len{P'}+1$.

Niech $Q$ będzie dowolną ścieżką z~$t$ do~$u$ w~$G$. Ponieważ $u\notin N[t]$, ścieżka $Q$ zaczyna się od krawędzi $ty$ z
$y\in N(t)$, więc jej obraz w~$G'$ ma długość nie większą niż $\len{Q}-1$. Stąd $\dist{G'}{t'}{u}\le \len{Q}-1$, czyli
$\len{Q}\ge \dist{G'}{t'}{u}+1=\len{P'}+1=\len{P}$. Otrzymujemy $\dist{G}{t}{u}=\len{P}$, a~zatem $P$ jest ścieżką najkrótszą
z~$t$ do~$u$ w~$G$.

Z~Obserwacji~\ref{obs:shortest-paths} ścieżka $P$ zawiera pewną krawędź $e\in F$. Pierwsza krawędź $tv$ ścieżki $P$ jest
incydentna z~$t$, więc nie należy do $F$ (bo $t\notin V[F]$). Zatem $e$ leży na odcinku $(v,x_1,\dots,x_k)$, który po ściągnięciu
$N[t]$ przechodzi w~odcinek $(t',x_1,\dots,x_k)$ ścieżki $P'$. W szczególności $e/N[t]$ jest krawędzią ścieżki $P'$ i~należy do $F'$,
więc $E(P')\cap F'\neq\emptyset$.

Jeżeli natomiast $u\neq t'$, $w\neq t'$ oraz $t'$ jest wierzchołkiem wewnętrznym ścieżki $P'$, to $P'$ dzieli się na konkatenację
$P'=P'_1\cdot P'_2$, gdzie $P'_1$ łączy $u$ z~$t'$, a~$P'_2$ łączy $t'$ z~$w$. Z~Lematu~\ref{lem:subpath} wynika, że obie podścieżki
są najkrótsze między swoimi końcami. Zastosowanie poprzedniego akapitu do par $(u,t')$ i~$(t',w)$ daje $E(P'_1)\cap F'\neq\emptyset$ oraz
$E(P'_2)\cap F'\neq\emptyset$, a~więc $E(P')\cap F'\neq\emptyset$.
\end{proof}

\section{Główne twierdzenie}

\begin{theorem}\label{thm:main}
Niech $G=(V,E)$ będzie grafem, $T\subseteq V$ oraz $F\subseteq E$. Jeżeli $F$ jest zbiorem semi-rozcinającym dla~$T$ w~$G$, to
\[
  |F|\ge |T|-1.
\]
\end{theorem}

\begin{proof}
Dowodzimy tezy przez indukcję po parze $(|T|,|V|)$ w~porządku leksykograficznym.

\medskip\noindent
\textbf{Baza indukcji.} Jeżeli $|T|\le 1$, to w~$T$ nie ma par różnych wierzchołków, więc warunek z~Definicji~\ref{def:semi-cut} jest
spełniony dla dowolnego $F$. W~takim przypadku $|F|\ge 0\ge |T|-1$.

\medskip\noindent
\textbf{Krok indukcyjny.} Niech $|T|\ge 2$ i~załóżmy, że twierdzenie jest prawdziwe dla wszystkich trójek $(G_0,T_0,F_0)$ z~parą
$(|T_0|,|V(G_0)|)$ mniejszą leksykograficznie od $(|T|,|V|)$. Rozróżnimy dwa przypadki.

\medskip\noindent
\textbf{Przypadek 1: $T\cap V[F]=\emptyset$.}
Wybierzmy dowolny $t\in T$ oraz dowolny $u\in T\setminus\{t\}$. Z~Definicji~\ref{def:semi-cut} wynika
$\dist{G}{t}{u}<\dist{G\setminus F}{t}{u}\le\infty$, a~zatem $\dist{G}{t}{u}<\infty$. W~szczególności istnieje ścieżka z~$t$ do~$u$,
więc $t$ ma sąsiada i~$|N[t]|\ge 2$.

Niech $G'\coloneqq G/N[t]$, $T'\coloneqq (T\setminus\{t\})\cup\{t'\}$ oraz $F'\coloneqq F/N[t]$ jak w~Propozycji~\ref{prop:contraction}.
Na mocy tej propozycji $F'$ jest zbiorem semi-rozcinającym dla~$T'$ w~$G'$. Ponadto $|T'|=|T|$ oraz, ponieważ $|N[t]|\ge 2$,
zachodzi $|V(G')|=|V|-|N[t]|+1<|V|$. Zatem z~założenia indukcyjnego otrzymujemy
\[
  |F'|\ge |T'|-1=|T|-1.
\]
Z~definicji obrazu zbioru krawędzi przy ściąganiu mamy $|F'|\le |F|$, a~więc $|F|\ge |T|-1$.

\medskip\noindent
\textbf{Przypadek 2: $T\cap V[F]\neq\emptyset$.}
Wybierzmy $t\in T\cap V[F]$ oraz krawędź $e\in F$ incydentną z~$t$. Zdefiniujmy $T^-\coloneqq T\setminus\{t\}$ oraz
$F^-\coloneqq F\setminus\{e\}$. Pokażemy, że $F^-$ jest zbiorem semi-rozcinającym dla~$T^-$ w~$G$.

Weźmy różne $u,w\in T^-$ i~niech $P$ będzie ścieżką z~$u$ do~$w$ o~długości $\len{P}=\dist{G}{u}{w}$.
Z~Obserwacji~\ref{obs:shortest-paths} wynika, że $P$ zawiera pewną krawędź ze zbioru $F$.
Jeżeli $e\notin E(P)$, to mamy od razu $E(P)\cap F^-\neq\emptyset$. Załóżmy więc, że $e\in E(P)$. Wtedy $t$ leży na~$P$ i~ścieżka
$P$ rozkłada się na konkatenację $P=P_u\cdot P_w$, gdzie $P_u$ łączy $u$ z~$t$, a~$P_w$ łączy $t$ z~$w$. Z~Lematu~\ref{lem:subpath}
wynika, że $\len{P_u}=\dist{G}{u}{t}$ oraz $\len{P_w}=\dist{G}{t}{w}$. Krawędź $e$ jest incydentna z~$t$, więc należy dokładnie do jednej z
podścieżek $P_u$ i~$P_w$. Jeżeli $e\notin E(P_w)$, to Obserwacja~\ref{obs:shortest-paths} zastosowana do pary $(t,w)$ daje
$E(P_w)\cap F\neq\emptyset$, a~zatem $E(P_w)\cap F^-\neq\emptyset$, co implikuje $E(P)\cap F^-\neq\emptyset$. Jeżeli natomiast
$e\notin E(P_u)$, to analogicznie z~pary $(u,t)$ otrzymujemy $E(P_u)\cap F^-\neq\emptyset$, więc $E(P)\cap F^-\neq\emptyset$.

Ponieważ $u,w$ były dowolne, z~Obserwacji~\ref{obs:shortest-paths} wynika, że $F^-$ jest zbiorem semi-rozcinającym dla~$T^-$ w~$G$.
Teraz para $(|T^-|,|V|)=(|T|-1,|V|)$ jest mniejsza leksykograficznie niż $(|T|,|V|)$, więc z~założenia indukcyjnego
\[
  |F^-|\ge |T^-|-1=|T|-2.
\]
W konsekwencji $|F|=|F^-|+1\ge |T|-1$, co kończy dowód.
\end{proof}

\end{document}
